\subsection{Experimental Results}
\vspace{-3pt}
As shown in \cref{tab:re10k_multiviews} and \cref{tab:crossdata_acid_nview}, our method achieves competitive performance among uncalibrated baselines and remains competitive with pose-free and pose-required methods. When using a similar number of Gaussian primitives as the baselines, our approach achieves strong reconstruction performance. Notably, even when using only 10-28\% of the Gaussian primitives, our method still maintains competitive results. These results suggest that our explicit density control enables a compact yet high-fidelity 3D Gaussian representation, leading to improved efficiency without sacrificing reconstruction quality.

\cref{fig:qual} shows qualitative comparisons on RE10K under different input-view settings. Our method produces sharper structures and more faithful scene details compared to existing approaches. Notably, even with substantially fewer Gaussian primitives (24-29\%), our approach maintains high rendering quality while reducing blurring artifacts, resulting in visually more accurate renderings.

\cref{tab:nvs_2view_acid} demonstrates the novel view synthesis performance on the ACID dataset under the challenging two-view setting. Our method achieves superior performance compared to the uncalibrated baseline. Moreover, it remains competitive with both pose-required and pose-free approaches, including feed-forward 3DGS methods and other baselines, demonstrating robust performance even with only two sparse input views.